% ECONOMICS_PROBLEM: {"id": "E71-141", "title": "Expectations beyond rational expectations", "jel_code": "E71", "source_id": "OP-0233"}
% !TeX program = xelatex
\documentclass[11pt,a4paper]{article}
\usepackage[margin=23mm]{geometry}
\usepackage{fontspec}
\setmainfont{Latin Modern Roman}
\usepackage{amsmath,amssymb,mathtools}
\usepackage{xcolor,enumitem,booktabs,longtable,array,xurl,hyperref}
\definecolor{muted}{HTML}{53636C}
\hypersetup{colorlinks=true,urlcolor=blue,linkcolor=blue}
\setlength{\parindent}{0pt}
\setlength{\parskip}{5pt}
\newcommand{\R}{\mathbb R}
\newcommand{\E}{\mathbb E}
\newcommand{\N}{\mathbb N}
\newcommand{\ind}{\mathbf 1}
\newcommand{\Var}{\operatorname{Var}}
\newcommand{\Cov}{\operatorname{Cov}}
\newcommand{\supp}{\operatorname{supp}}
\DeclareMathOperator*{\argmin}{arg\,min}
\DeclareMathOperator*{\argmax}{arg\,max}
\newcommand{\entrymeta}[3]{{\sffamily\small\color{muted}\textbf{\texttt{#1}}\quad|\quad #2\par #3\par}}
\newcommand{\Problemref}[1]{\texttt{#1}}
\newcommand{\JELref}[2]{\texttt{#2} (JEL \texttt{#1})}
\begin{document}
\section*{Expectations beyond rational expectations}
\textbf{Problem ID:} \texttt{E71-141}\quad \textbf{JEL:} \texttt{E71}\par
% BEGIN ECONOMICS STATEMENT
\label{problem:OP-0233}
\label{audited:ECON-027}
{\sffamily\small\color{muted}Original subject group: Business cycles and macro dynamics\par}
\label{definitions:problem:30007514}
\textbf{Audit classification:} Explicit published open question. The 2026-displayed manuscript explicitly proposes three criteria; the architecture and accuracy targets are editorial. Verification concerns the cited version, not first priority or proof that this exact editorial specification remains unsolved on October 8, 2026.
\subsection*{Definitions and precise research target}
\paragraph{Editorial mathematical specification.} Fix an incomplete-markets economy with heterogeneous households, idiosyncratic income risk, aggregate shocks, and a borrowing constraint. Let \(\mu_t\) denote the joint distribution of assets and income, \(s_t\) the aggregate fundamental state, and \(p_t\) the vector of equilibrium prices. Each household forecasts future prices through a subjective conditional distribution \(Q_\theta(\,\cdot\mid I_{it},b_{it})\), where \(I_{it}\) is its observed information and \(b_{it}\) is a finite-dimensional belief state. Specify a updating rule \(b_{i,t+1}=G_\theta(b_{it},I_{i,t+1},a_t)\), with \(a_t\) the observed policy regime. Expectations in household optimality conditions are taken under \(Q_\theta\); market clearing determines actual prices.
Construct and identify a pair \((Q_\theta,G_\theta)\) satisfying three jointly assessed requirements: computationally feasible global equilibrium with aggregate risk; agreement with held-out household forecast distributions and errors; and endogenous, testable changes in beliefs following policy changes. More precisely, fix tolerances \(\epsilon_F,\epsilon_M>0\) and require forecast-distance and macroeconomic prediction loss to fall below them on a declared evaluation sample. Demonstrate which restrictions identify \(\theta\) and prevent arbitrary fitting through unobserved beliefs. The target is a disciplined alternative to full rational expectations in this economy, with reproducible validation and sensitivity analysis, rather than proof that any one learning rule universally replaces rational expectations.

Use a probability space carrying actual fundamental shocks and a filtration \(\mathcal F_t\); household information \(\mathcal I_{it}\subseteq\mathcal F_t\) generates \(I_{it}\). Each \(Q_\theta\) must be a probability kernel on future price paths with measurable dependence on information, and \(G_\theta\) a measurable belief transition. For example, household feasibility is \(c_{it}+a_{i,t+1}=(1+r_t)a_{it}+w_te_{it}\ell_{it}+T_{it}\), with \(a_{i,t+1}\ge\underline a\), \(c_{it}>0\), and a specified utility. A temporary equilibrium maximizes subjective expected utility and clears goods, assets, and labor, while actual distributions evolve under the actual shock law. Specify a computational budget, a proper forecast score, and bounded macro loss before choosing \(\theta\); require the empirical feasible set of pairs meeting both tolerances to be nonempty, or report failure rather than assuming success.
\subsection*{Variants and assumption changes}
The following are \emph{editorial variants}, not additional published open conjectures.
\begin{enumerate}
\item \emph{Common adaptive beliefs.} Impose a common finite-order least-squares price forecast and a specified gain sequence. Study whether its stationary learning equilibrium exists and whether its out-of-sample forecast score meets the declared tolerances.
\item \emph{Heterogeneous information.} Let households observe different signals and carry different belief states. Hold signal precision fixed across policies and require the same updating parameters to predict responses to new policy announcements; compare with the common-belief restriction.
\end{enumerate}
\subsection*{References and source audit}
\begin{itemize}
\item Benjamin Moll. \emph{Research Agenda: Benjamin Moll}. 2020. Locator: Section4.1. Primary text checked. \url{https://economicdynamics.org/research-agenda-moll2020/}.
\item Benjamin Moll. \emph{The Trouble with Rational Expectations in Heterogeneous Agent Models: A Challenge for Macroeconomics}. 2025. Locator: Section3.3 (three criteria); displayed manuscript version August 24, 2026; arXiv banner dated August 28, 2025. Primary text checked. \url{https://arxiv.org/html/2508.20571v1}.
\end{itemize}
The 2026-displayed manuscript explicitly proposes three criteria; the architecture and accuracy targets are editorial.
\begin{enumerate}\raggedright
\item\hypertarget{definitions:source:30007514:1}{} Benjamin Moll. 2020. \emph{Research Agenda: Benjamin Moll}. Section4.1.
\url{https://economicdynamics.org/research-agenda-moll2020/}.
\item\hypertarget{definitions:source:30007514:2}{} Benjamin Moll. 2025. \emph{The Trouble with Rational Expectations in Heterogeneous Agent Models: A Challenge for Macroeconomics}. Section3.3 (three criteria); displayed manuscript version August24,2026; arXiv banner dated August28,2025.
\url{https://arxiv.org/html/2508.20571v1}.
DOI: \url{https://doi.org/10.48550/arXiv.2508.20571}.
\end{enumerate}

\subsection*{Integrated refinement: ECON-027}
{\sffamily\small\color{muted}Empirically disciplined heterogeneous beliefs in macroeconomics\par}
This refinement adopts the additional source conventions
(page~\pageref{audited:conventions}), subject to its local restrictions.
\textbf{Replacing heterogeneous beliefs by an objective-law benchmark.}
For exogenous states with realized law $P$, subjective kernels $Q_i$ and belief
distribution $F_Q$, define a selected-equilibrium aggregate $Y(F_Q)$ and
\[
\Delta Y=Y(F_Q)-Y(\delta_P).
\]
Use joint survey forecasts and choices with an explicit measurement-error model
to identify the belief distribution sufficiently for this contrast; otherwise
report its identified set. Both economies have the same realized exogenous law
and initial resources. Endogenous prices and their forecast laws must be solved
again: replacing beliefs by the old equilibrium price law need not give a
self-consistent rational-expectations benchmark. Distinguish disagreement from
private information and preferences, learning from fixed beliefs, and
fixed-price comparisons from a new equilibrium.
\subsection*{Supplied source audit for ECON-027}
Local [S1], [S2], etc. in this audit and the references immediately below
belong to ECON-027, independently of the original entry's reference numbering.
[S1], PDF p. 6, lists heterogeneous beliefs and empirical discipline. This counterfactual is editorial and does not claim that survey forecasts uniquely identify macroeconomic beliefs.
\subsection*{References for integrated refinement ECON-027}
{\footnotesize\raggedright
\par\noindent\textbf{[S1]} Benjamin Moll (n.d.). \emph{Lecture 11: Good Research Topics and Open Questions}. Distributional Macroeconomics lecture slides; publicly archived at a 2019 URL.\par \textit{Verified locator / source role:} PDF p. 6: heterogeneous beliefs and empirical discipline.\par \url{https://benjaminmoll.com/wp-content/uploads/2019/07/Lecture11_2149.pdf}\par\smallskip
}

\subsection*{Common conventions: Source mathematical conventions}

These conventions apply to each entry unless explicitly replaced.

\textbf{Models and identification.} A primitive model \(m\) specifies all probability laws, feasible choices, information, equilibrium and measurement rules used by its target. The admissible class \(\mathcal M\) is fixed independently of outcomes. If \(P_O(m)\) is its observable probability law and \(\tau(m)\) the target, its identified set at observable law \(P\) is
\[
 \mathcal I_\tau(P)=\{\tau(m):m\in\mathcal M,\ P_O(m)=P\}.
\]
Point identification means that this set is a singleton. An empty set signals incompatibility of data and assumptions. Coordinate bounds need not describe the joint set. Bounds are sharp only if every retained target is attainable or arbitrarily approachable by compatible models. Finite bounds require explicit restrictions on unobserved quantities; unrestricted confounding, hidden wealth, extrapolation or arbitrary equilibrium selection can yield uninformative sets.

\textbf{Causal interventions.} A potential outcome \(Y_i(p)\) is the outcome under a complete policy regime \(p\), including timing, financing, information and market spillovers. The target population and units are fixed before treatment. With interference, use the complete assignment vector or a justified exposure mapping; individual no-interference notation alone cannot identify an economy-wide intervention. A difference of mean potential outcomes does not require the cross-world joint distribution. Conditioning on post-treatment migration, survival, enrollment or employment changes the estimand and needs separate justification.

\textbf{Statistical statements.} An estimator, confidence set, forecast or test specifies a sampling law, model class and loss/coverage criterion. Uniform \((1-\alpha)\)-coverage means \(\inf_{m\in\mathcal M}\Pr_m\{\tau(m)\in C_n\}\geq1-\alpha\), or an explicitly asymptotic version. The whole parameter space trivially has coverage, so informativeness requires a stated width, risk or power objective. A forecasting improvement is relative to a named benchmark, information set and evaluation population.

\textbf{Equilibrium and optimization.} An equilibrium comprises feasible plans, optimality, beliefs where needed, budget constraints and all market/resource clearing. The primitives or a stated benchmark define each correspondence \(\mathcal E(p;m)\). Nonemptiness is assumed or proved, never inferred just from compact policy sets. A selected-equilibrium target uses a declared selection rule; a robust target quantifies over all permitted equilibria. Use suprema or infima unless attainment is established. An \(\varepsilon\)-optimal policy is feasible and within \(\varepsilon\) of the finite optimum value. Report an empty feasible set separately.

\textbf{Welfare and accounting.} Normative utility, interpersonal normalization, social weights, numeraire, discounting and the use of public receipts are primitives. Behavioral utility need not coincide with the welfare criterion. Household welfare includes ownership income and rents once; adding profits again to compensating variation generally double-counts them. A monetary equivalent is defined by an explicit transfer and price convention, with monotonicity and range conditions. Average effects, mechanism contributions, transfers and welfare losses are different targets. Infinite sums require convergence and dynamic feasible plans require terminal or no-Ponzi conditions.

\textbf{Source versions.} A working-paper date, revision date and journal publication year are distinguished. A DOI for a journal version is not automatically the DOI of an earlier manuscript. Locators refer to the stated version and printed pages where specified. A metadata-only check does not verify a claim about a full-text conclusion. Every entry's source subsection narrows the provenance claim accordingly.

\subsection*{Common conventions: Additional-source status and mathematical conventions}

\label{audited:conventions}

Of the 100 supplied ECON entries, 65 distinct problems appear here; 32 close
matches refine earlier entries, and three duplicate questions map to existing
entries. Appendix~\ref{app:audited-crosswalk} lists every destination.
The attachment's P/R/S/U distinctions and source audits are retained. Its
precise mathematical formalizations are editorial unless the individual audit
says otherwise. Candidate subsidy targets ECON-004--006 retain their U flags;
the resolved approval-core question ECON-007 updates OP-0202 above.
The following supplied conventions govern both this part and the attributed
ECON refinements elsewhere in the catalogue, subject to local restrictions.

\textbf{Parameterized questions.} A problem family lists inputs that must be fixed before an instance is evaluated: population, horizon, policies, information, data law, model class and welfare weights. These are required choices, not quantities the citations are assumed to specify. Undefined applications should not be treated as identified estimands.

\textbf{Indices and integrability.} Sets of agents, firms and regions are finite unless a population measure is explicitly used. \(i,j\) index units, \(r\) regions, \(t\) dates and \(h\) horizons. \(\mathbb E\) and \(\Pr\) refer to the declared population and law; \(\mathbf1\{A\}\) is an indicator. Every displayed expectation and discounted sum needs existence in the stated domain. Infinite horizons use a discount in \((0,1)\) and a convergence condition; finite horizons may use discount one.

\textbf{Optimization and existence.} A displayed \(\max\), \(\min\), or \(\arg\max\) is conditional on attainment. For finite feasible menus this is immediate; general spaces need continuity and compactness/coercivity or another existence argument. Without attainment, the target is the supremum/infimum and arbitrarily near-optimal feasible choices. \(\inf\varnothing=+\infty\). Empty feasibility and an unidentified target are different issues.

\textbf{Potential outcomes.} \(Y_i(z)\) is the outcome under a specified intervention, not the observed conditional mean among people choosing \(z\). In binary comparisons \(\Delta Y_i=Y_i(1)-Y_i(0)\). Specify assignment, treatment versions, compliance, information and observation horizon. Interference is allowed under a full policy vector; compressed exposure notation needs a justified mapping. No interference, consistency and overlap are substantive assumptions, not automatic consequences of notation.

\textbf{Populations and observation.} Fix baseline eligibility and population weights. Conditioning on post-policy employment, survival, relocation or program uptake can change the population being compared. Death is not ordinary loss to follow-up; outcomes truncated by death need a composite convention, joint survival target, or explicitly defined survivor stratum. Fertility-changing interventions need a population functional rather than an unexplained fixed descendant cohort. Measurement and missingness models must be supplied.

\textbf{Identification.} A structural model \(m\in\mathcal M\) implies observable law \(P_m\) and target \(\theta(m)\). Identification means
\[
 P_{m_1}=P_{m_2}\ \Longrightarrow\ \theta(m_1)=\theta(m_2)
 \quad\text{for all }m_1,m_2\in\mathcal M .
\]
Without identification, the target is
\[
 \Theta(P;\mathcal M)=\{\theta(m):m\in\mathcal M,\ P_m=P\}.
\]
Writing \(\theta(P)\) before establishing identification can hide incompatible latent models. Confidence sets additionally need a sampling law, confidence level, asymptotic sequence and uniformity class. Point identification, coverage and informative precision are separate properties.

\textbf{Mediation.} Total effects of randomized assignment are distinct from cross-world natural indirect effects. Belief/effort-channel effects require a meaningful mediator intervention and additional measurement, exclusion or mediation assumptions. Randomization of the initial policy alone does not supply those assumptions. Report bounds or interventional alternatives when the stronger target is unsupported.

\textbf{Equilibrium.} A counterfactual \(W(z)\), \(p(z)\), or \(\mu_t^z\) requires individual optimization, feasibility, government/resource budgets, market clearing and state-transition laws. State equilibrium existence and selection, or report ranges over compatible equilibria. Initial-state integration includes subsequent endogenous transitions; current-state integration must use the policy-dependent distribution. Reduced-form invariance after policy is not assumed.

\textbf{Welfare accounts.} Scores, utility, health, dollars and ecological quantities require explicit conversion before addition. Fees, taxes, subsidies, wages and extortion payments are transfers between parties; distributional weights can make them welfare relevant, but their face value is not automatically a resource loss. State ownership, geographic boundaries, foreign-income weights and fiscal finance. Do not add profits to household welfare already including dividends, or count land capitalization and the underlying benefit twice. A benefit-cost expression must distinguish gross benefits from net welfare and subtract every real cost exactly once.

\textbf{Derivatives and risk.} Log derivatives require positive arguments; local responses require admissible perturbations and specified equilibrium branches. Differentiation under expectations needs regularity/domination, and boundaries may allow only one-sided derivatives. Risk uses a specified law; ambiguity uses a declared nonempty law set on a common history space. Adaptive policies must be nonanticipative. A robust criterion is an editorial choice unless attributed explicitly.

\textbf{Scope overlaps.} Allocation, collective choice, contracts, household bargaining and coordinated policy overlap economics with optimization and game theory. They are retained because they appeared in the original catalogue; no claim of a strictly game-theory-free selection is made.
% END ECONOMICS STATEMENT
\end{document}
